% STATUS: DRAFTED 2026-09-21 by the Search Process Policy agent. Full section
% written per D10 + Section Blueprint. Citations added to
% iclr2026_conference_policy_refs.bib (hart1968formal, pohl1970heuristic,
% bonet2001planning, lipovetzky2017bfws, ross2011dagger; all verified against
% Crossref/PMLR metadata on 2026-09-21). No new macros; only frozen labels
% sec:design and appendix are referenced.
\section{Search Process Policy}
\label{sec:policy}

% Paragraph purpose: Position the section: state the formal contract behind the mechanism introduced in Section 1, with pointers to protocol and appendix material.
The introduction described the Search Process Policy mechanism informally. This section states the execution contract it runs under: the task inputs, the Modality Observation presented to the policy, the Typed Search Operation envelope, the Search Memory held by the runtime, the Algorithm Invariants checked on every operation, the ownership rule separating model decisions from runtime bookkeeping, and the learning objective. Training conditions, panels, and budgets are protocol decisions and appear in Section~\ref{sec:design}. Concrete operation schemas, invariant definitions, and runtime pseudocode appear in Appendix~\ref{appendix}.

% Paragraph purpose: Define the task instance and episode configuration formally and fix the active algorithm set.
A task instance is a triple $x = (D, s_0, G)$, where $D$ is a formal domain model whose grounded actions induce a finite state space with deterministic transitions, $s_0$ is the initial state, and $G$ is a goal condition over states. An episode fixes a tuple $(x, A, m, B)$: an instance $x$, a declared search algorithm $A$, a modality adapter $m$, and a frozen budget $B$ bounding the operations the episode may issue. The declared algorithm is part of the input, not a choice made at run time. The active algorithm set contains breadth-first search (BFS), the complete unpruned best-first width search (BFWS) with novelty and goal-count priorities \citep{lipovetzky2017bestfirst}, and two additive best-first variants over the $h_{\rm add}$ heuristic \citep{bonet2001planning}: weighted A* with weight $w = 3$ \citep{hart1968formal,pohl1970heuristic} and greedy best-first.

% Paragraph purpose: Define the Modality Observation and its fixed adapter contract, including decision sufficiency.
The policy never reads the instance directly. At decision step $t$ it receives a Modality Observation $o_t = \phi_m(c_t)$, where the decision context $c_t$ contains the canonical task context and the decision-relevant view of the current Search Memory, and $\phi_m$ is a fixed observation adapter. Three adapters are declared. text-state serializes the state, the goal, and the search context as a compact relational representation. visual-state renders the same state and goal as a rendered state image plus a partial-goal constraint image. multimodal-state presents both. In all three adapters, Search Memory contents and the grounded candidate list remain in the textual serialization. The adapters differ only in how the state and goal are presented. The adapter is fixed for an entire training or evaluation cell, and its contract requires a decision-sufficient observation: $o_t$ must expose the facts the exact teacher uses to resolve step $t$, such as the grounded successor candidates and visited membership, so that a correct operation is reachable from the observation alone.

% Paragraph purpose: Define the Typed Search Operation envelope and the operand ownership rule.
The policy output at step $t$ is a single Typed Search Operation $u_t = (\tau_t, \theta_t)$. The operation type $\tau_t$ comes from the vocabulary of the declared algorithm, covering steps such as expanding a node, generating a successor, updating the frontier, or applying a goal test. The operand vector $\theta_t$ carries the arguments of that step. Ownership is defined over operands. Every operand that determines exploration must be emitted by the policy, including the selected source or frontier state and any successor, action, or insertion decision the declared operation requires. An emission that omits an exploration-determining operand, or names operands inconsistent with the decision context, is invalid by definition.

% Paragraph purpose: Define Search Memory as external, bounded, runtime-held state distinct from model state.
Search Memory $M_t$ is the external state the Trusted Search Runtime maintains across an episode. It holds the frontier, the visited or best-depth table, and algorithm-specific state such as novelty counters for BFWS. It is bounded, declared per protocol cell, and it is not internal unbounded model state.
% TODO: the fixed Search Memory capacity value is an open preregistration
% prerequisite (writing_design_tree.md); pin the value here once frozen.
The policy has no access to episode history except through what the adapter exposes from $M_t$ in $o_t$. Equivalently, the observation sequence is the only channel through which prior decisions influence later ones.

% Paragraph purpose: Define Algorithm Invariants per algorithm and delimit what stepwise validity does and does not establish.
Each declared algorithm carries Algorithm Invariants, deterministic properties the runtime checks on every operation. For BFS, the frontier must be consumed in first-in-first-out order. For BFWS, selection must follow novelty and goal-count priority with the declared duplicate handling. For the additive best-first variants, frontier order must follow $f(s) = g(s) + w \cdot h_{\rm add}(s)$ with $w = 3$, where $g(s)$ is the accumulated cost of the path that reached $s$, or $h_{\rm add}$ alone for greedy best-first, under the declared tie-breaking rule.
% TODO: the per-algorithm tie-breaking convention must be recorded in the
% Appendix evidence contract; the exact rule is not pinned in current sources.
Stepwise validity and invariant compliance are execution-quality measurements. They do not establish termination, completeness, optimality, or episode success, which are reported separately under each cell's declared task, heuristic, tie-breaking rule, and budget.

% Paragraph purpose: State the model/runtime ownership rule and the replay-determinative attribution test.
The Trusted Search Runtime has four permitted actions on an emission: validate it, apply it, persist its effects to Search Memory, or reject it. It may not choose omitted operands, reorder candidates on the policy's behalf, silently repair an invalid emission, or substitute a default search decision. This boundary gives the attribution test used throughout the paper. Replaying the raw policy emissions $u_{1:T}$ against the same task and the same initial Search Memory must reproduce the explored trace exactly. A decision that appears only because the runtime filled, reordered, or repaired an emission is runtime behavior, not model behavior, and the replay test detects it.

% Paragraph purpose: Give one generic worked micro-example of a single step under BFS.
A single step under BFS illustrates the contract. The observation lists the frontier in order and, for the current source state, the grounded successor candidates with their visited membership. To expand, the policy emits an operation whose operands name the source state and, for each insertion, the chosen successor and its enqueue position. The runtime validates that the named source is the frontier head, that each named successor is an unvisited grounded successor of that source, and that each insertion lands at the tail, then applies and persists the update. Any failed check rejects the emission. Nothing is repaired, and no alternative source, successor, or position is chosen on the policy's behalf.

% Paragraph purpose: Define the learning objective: process supervision on trusted Search-Trace Segments, with optional DAgger iterations.
The training unit is a Search-Trace Segment, a bounded ordered slice of Typed Search Operations and the observations that preceded them, cut from replay-verified traces of exact algorithm execution. Process supervision is behavior cloning over a dataset $\train$ of such segments. The policy is trained with token-level cross-entropy to emit the teacher's operation given the observation, so supervision lands on the exploration decisions themselves rather than on a final plan. An optional DAgger stage \citep{ross2011dagger} rolls the current policy out under the runtime, collects teacher corrections on the decision contexts the policy actually visits, aggregates them into $\train$, and retrains. Both conditions train against the same contract above. The arms, exposure matching, and evaluation gates that compare them are defined in Section~\ref{sec:design}.
